\documentclass[10pt,letterpaper]{article}
\usepackage[utf8]{inputenc}
\usepackage[intlimits]{amsmath}
\usepackage{amsfonts}
\usepackage{amssymb}
\usepackage{ragged2e}
\usepackage[letterpaper, margin=0.5in]{geometry}
\usepackage{graphicx}
\usepackage{cancel}
\usepackage{mathtools}
\usepackage{tabularx}
\usepackage{arydshln}
\usepackage{tensor}
\usepackage{array}
\usepackage{xcolor}
\usepackage[boxed]{algorithm}
\usepackage[noend]{algpseudocode}
\usepackage{listings}
\usepackage{textcomp}
% \usepackage[pdf,tmpdir,singlefile]{graphviz}
\usepackage{mathrsfs}
\usepackage{bbm}
\usepackage{tikz}
\usepackage{tikz-cd}
\usepackage{enumitem}
\usepackage{arydshln}
\usepackage{relsize}
\usepackage{multirow,multicol}
\usepackage{scalerel}
\usepackage{upgreek}
\usepackage{ifthen}
\usepackage{yhmath}

\usetikzlibrary{bayesnet}
\usetikzlibrary{patterns}
\usetikzlibrary{calc}
\setlist{noitemsep}

%%%%%%%%%%%%%%%%%%%%%%%%%%%%%
% Formatting commands
%%%%%%%%%%%%%%%%%%%%%%%%%%%%%
\newcommand{\n}{\hfill\break}
\newcommand{\nn}{\vspace{0.5\baselineskip}\n}
\newcommand{\up}{\vspace{-\baselineskip}}
\newcommand{\hangblock}[2]{\par\noindent\settowidth{\hangindent}{\textbf{#1: }}\textbf{#1: }\nolinebreak#2}
\newcommand{\lemma}[1]{\hangblock{Lemma}{#1}}
\newcommand{\defn}[1]{\hangblock{Defn}{#1}}
\newcommand{\thm}[1]{\hangblock{Thm}{#1}}
\newcommand{\cor}[1]{\hangblock{Cor}{#1}}
\newcommand{\prop}[1]{\hangblock{Prop}{#1}}
\newcommand{\ex}[1]{\hangblock{Ex}{#1}}
\newcommand{\exer}[1]{\hangblock{Exer}{#1}}
\newcommand{\fact}[1]{\hangblock{Fact}{#1}}
\newcommand{\remark}[1]{\hangblock{Remark}{#1}}
\newcommand{\proven}{\;$\square$\n}
\newcommand{\problem}[1]{\par\noindent{\nolinebreak#1}\n}
\newcommand{\problempart}[2]{\par\noindent\indent{}\settowidth{\hangindent}{\textbf{(#1)} \indent{}}\textbf{(#1) }\nolinebreak#2\n}
\newcommand{\ptxt}[1]{\textrm{\textnormal{#1}}}
\newcommand{\inlineeq}[1]{\centerline{$\displaystyle #1$}}
\newcommand{\pageline}{\up\par\noindent\rule{\textwidth}{0.1pt}}

%%%%%%%%%%%%%%%%%%%%%%%%%%%%%
% Math commands
%%%%%%%%%%%%%%%%%%%%%%%%%%%%%
% Set Theory
\newcommand{\card}[1]{\left|#1\right|}
\newcommand{\set}[1]{\left\{#1\right\}}
\newcommand{\setmid}{\;\middle|\;}
\newcommand{\ps}[1]{\mathcal{P}\left(#1\right)}
\newcommand{\pfinite}[1]{\mathcal{P}^{\ptxt{finite}}\left(#1\right)}
\newcommand{\naturals}{\mathbb{N}}
\newcommand{\N}{\naturals}
\newcommand{\integers}{\mathbb{Z}}
\newcommand{\Z}{\integers}
\newcommand{\rationals}{\mathbb{Q}}
\newcommand{\Q}{\rationals}
\newcommand{\reals}{\mathbb{R}}
\newcommand{\R}{\reals}
\newcommand{\complex}{\mathbb{C}}
\newcommand{\C}{\complex}
\newcommand{\halfPlane}{\mathbb{H}}
\let\HSym\H
\let\H\relax
\newcommand{\H}{\halfPlane}
\newcommand{\comp}{^{\complement}}
\DeclareMathOperator{\Hom}{Hom}
\newcommand{\Ind}{\mathbbm{1}}
\newcommand{\cut}{\setminus}
\DeclareMathOperator{\elem}{elem}

% Differentiable Manifolds
\newcommand{\RP}{\mathbb{RP}}
\newcommand{\CP}{\mathbb{RP}}
\newcommand{\osubset}{\overset{\mathclap{\scalebox{0.5}{\ptxt{open}}}}{\subset}}
\newcommand{\osubseteq}{\overset{\mathclap{\scalebox{0.5}{\ptxt{open}}}}{\subseteq}}
\newcommand{\osupset}{\overset{\mathclap{\scalebox{0.5}{\ptxt{open}}}}{\supset}}
\newcommand{\osupseteq}{\overset{\mathclap{\scalebox{0.5}{\ptxt{open}}}}{\supseteq}}
\newcommand{\pdat}[3]{\left.\pd{#1}{#2}\right|_{#3}}
\DeclareMathOperator{\so}{so}
\DeclareMathOperator{\codim}{codim}
\DeclareMathOperator{\Diff}{Diff}
\let\dSym\d
\let\d\relax
\newcommand{\d}{\partial}
\DeclareMathOperator{\gl}{gl}
\DeclareMathOperator{\Ad}{Ad}
\newcommand{\comm}[1]{\left[#1\right]}
\DeclareMathOperator{\Fix}{Fix}
\DeclareMathOperator{\Gr}{Gr}
\DeclareMathOperator{\Ric}{Ric}
\let\OSymbol\O
\let\O\relax
\DeclareMathOperator{\O}{O}
\newcommand{\overarc}[1]{\wideparen{#1}}
\DeclareMathOperator{\II}{\mathit{I\!\!I}}
\DeclareMathOperator{\Vol}{Vol}
\DeclareMathOperator{\Pf}{Pf}
\DeclareMathOperator{\Laplacian}{\Delta}

% Graph Theory
\let\deg\relax
\DeclareMathOperator{\deg}{deg}
\newcommand{\degp}{\ptxt{deg}^{+}}
\newcommand{\degn}{\ptxt{deg}^{-}}
\newcommand{\precdot}{\mathrel{\ooalign{$\prec$\cr\hidewidth\hbox{$\cdot\mkern0.5mu$}\cr}}}
\newcommand{\succdot}{\mathrel{\ooalign{$\cdot\mkern0.5mu$\cr\hidewidth\hbox{$\succ$}\cr\phantom{$\succ$}}}}
\DeclareMathOperator{\cl}{cl}
\DeclareMathOperator{\affdim}{affdim}

% Probability
\newcommand{\parSymbol}{\P}
\newcommand{\Prob}{\mathbb{P}}
\renewcommand{\P}{\Prob}
\newcommand{\Avg}{\mathbb{E}}
\newcommand{\E}{\Avg}
\DeclareMathOperator{\Var}{Var}
\DeclareMathOperator{\Cov}{Cov}
\DeclareMathOperator{\cov}{cov}
\DeclareMathOperator{\Corr}{Corr}
\DeclareMathOperator{\Unif}{Unif}
\DeclareMathOperator{\Binom}{Binom}
\newcommand{\CI}{\mathrel{\text{\scalebox{1.07}{$\perp\mkern-10mu\perp$}}}}
\DeclareMathOperator{\Ber}{Ber}
\DeclareMathOperator{\Bin}{Bin}
\DeclareMathOperator{\Geom}{Geom}
\DeclareMathOperator{\Poisson}{Poisson}
\DeclareMathOperator{\Exp}{Exp}
\DeclareMathOperator{\Beta}{Beta}
\DeclareMathOperator{\Normal}{N}
\DeclareMathOperator{\DistGamma}{Gamma}
\DeclareMathOperator{\eff}{eff}
\DeclareMathOperator{\AVar}{AVar}

% Standard Math
\newcommand{\inv}{^{-1}}
\newcommand{\abs}[1]{\left|#1\right|}
\newcommand{\ceil}[1]{\left\lceil{}#1\right\rceil{}}
\newcommand{\floor}[1]{\left\lfloor{}#1\right\rfloor{}}
\newcommand{\conj}[1]{\overline{#1}}
\newcommand{\of}{\circ}
\newcommand{\tri}{\triangle}
\newcommand{\inj}{\hookrightarrow}
\newcommand{\surj}{\twoheadrightarrow}
\newcommand{\ndiv}{\nmid}
\renewcommand{\epsilon}{\varepsilon}
\newcommand{\divides}{\mid}
\newcommand{\ndivides}{\nmid}
\DeclareMathOperator{\lcm}{lcm}
\DeclareMathOperator{\sgn}{sgn}
\newcommand{\map}[4]{{\arraycolsep=2pt\begin{array}[t]{rcl}#1 & \to & #2\\ {}#3 & \mapsto & #4\end{array}}}
\newcommand{\bigsum}[2]{\smashoperator[lr]{\sum_{\scalebox{#1}{$#2$}}}}
\DeclareMathOperator{\gcf}{gcf}
\newcommand{\restr}[1]{\left.#1\right|}

% Linear Algebra
\newcommand{\Id}{\textrm{\textnormal{Id}}}
\newcommand{\im}{\textrm{\textnormal{im}}}
\newcommand{\norm}[1]{\abs{\abs{#1}}}
\newcommand{\tpose}{^{T}\!}
\newcommand{\iprod}[1]{\left<#1\right>}
\newcommand{\giprod}{\iprod{\;\,,\;}}
\DeclareMathOperator{\tr}{tr}
\DeclareMathOperator{\trace}{tr}
\newcommand{\chgBasMat}[3]{\!\!\tensor*[_{#1}]{\left[#2\right]}{_{#3}}}
\newcommand{\vecBas}[2]{\tensor*[]{\left[#1\right]}{_{#2}}}
\DeclareMathOperator{\GL}{GL}
\DeclareMathOperator{\Mat}{Mat}
\DeclareMathOperator{\vspan}{span}
\DeclareMathOperator{\rank}{rank}
\newcommand{\V}[1]{\vec{#1}}
\DeclareMathOperator{\proj}{proj}
\DeclareMathOperator{\compProj}{comp}
\DeclareMathOperator{\row}{row}
\newcommand{\smallPMatrix}[1]{\paren{\begin{smallmatrix}#1\end{smallmatrix}}}
\newcommand{\smallBMatrix}[1]{\brack{\begin{smallmatrix}#1\end{smallmatrix}}}
\newcommand{\pmat}[1]{\begin{pmatrix}#1\end{pmatrix}}
\newcommand{\bmat}[1]{\begin{bmatrix}#1\end{bmatrix}}
\newcommand{\dual}{^{*}}
\newcommand{\pinv}{^{\dagger}}
\newcommand{\horizontalMatrixLine}{\ptxt{\rotatebox[origin=c]{-90}{$|$}}}
\DeclareMathOperator{\range}{range}
\DeclareMathOperator{\Symm}{Symm}
\DeclareMathOperator{\SU}{SU}
\DeclareMathOperator{\U}{U}

% Multilinear Algebra
\let\Lsym\L
\let\L\relax
\DeclareMathOperator{\L}{\mathscr{L}}
\DeclareMathOperator{\A}{\mathcal{A}}
\DeclareMathOperator{\Alt}{Alt}
\DeclareMathOperator{\Sym}{Sym}
\newcommand{\ot}{\otimes}
\newcommand{\ox}{\otimes}
\DeclareMathOperator{\asc}{asc}
\DeclareMathOperator{\asSet}{set}
\DeclareMathOperator{\sort}{sort}
\DeclareMathOperator{\ringA}{\mathring{A}}
\DeclareMathOperator{\Sh}{Sh}
\DeclareMathOperator{\Bil}{Bil}

% Topology
\newcommand{\closure}[1]{\overline{#1}}
\newcommand{\uball}{\mathcal{U}}
\DeclareMathOperator{\Int}{Int}
\DeclareMathOperator{\Ext}{Ext}
\DeclareMathOperator{\Bd}{Bd}
\DeclareMathOperator{\rInt}{rInt}
\DeclareMathOperator{\ch}{ch}
\DeclareMathOperator{\ah}{ah}
\newcommand{\LargerTau}{\mathlarger{\mathlarger{\mathlarger{\mathlarger{\tau}}}}}
\newcommand{\Tau}{\mathcal{T}}

% Analysis
\DeclareMathOperator{\Graph}{Graph}
\DeclareMathOperator{\epi}{epi}
\DeclareMathOperator{\hypo}{hypo}
\DeclareMathOperator{\supp}{supp}
\newcommand{\lint}[2]{\underset{#1}{\overset{#2}{{\color{black}\underline{{\color{white}\overline{{\color{black}\int}}\color{black}}}}}}}
\newcommand{\uint}[2]{\underset{#1}{\overset{#2}{{\color{white}\underline{{\color{black}\overline{{\color{black}\int}}\color{black}}}}}}}
\newcommand{\alignint}[2]{\underset{#1}{\overset{#2}{{\color{white}\underline{{\color{white}\overline{{\color{black}\int}}\color{black}}}}}}}
\newcommand{\extint}{\ptxt{ext}\int}
\newcommand{\extalignint}[2]{\ptxt{ext}\alignint{#1}{#2}}
\newcommand{\conv}{\ast}
\newcommand{\pd}[2]{\frac{\partial{}#1}{\partial{}#2}}
\newcommand{\del}{\nabla}
\DeclareMathOperator{\grad}{grad}
\DeclareMathOperator{\curl}{curl}
\let\div\relax
\DeclareMathOperator{\div}{div}
\DeclareMathOperator{\vol}{vol}

% Complex Analysis
\let\Re\relax
\DeclareMathOperator{\Re}{Re}
\let\Im\relax
\DeclareMathOperator{\Im}{Im}
\DeclareMathOperator{\Res}{Res}

% Abstract Algebra
\DeclareMathOperator{\ord}{ord}
\newcommand{\generated}[1]{\left<#1\right>}
\newcommand{\cycle}[1]{\smallPMatrix{#1}}
\newcommand{\id}{\ptxt{id}}
\newcommand{\iso}{\cong}
\DeclareMathOperator{\Aut}{Aut}
\DeclareMathOperator{\SL}{SL}
\DeclareMathOperator{\op}{op}
\newcommand{\isom}[4]{\!\!\!\begin{array}[t]{rcl}#1 & \!\!\!\!\overset{\sim}{\to} & \!\!\!\!#2\\ #3 & \!\!\!\!\mapsto & \!\!\!\!#4\end{array}}
\newcommand{\F}{\mathbb{F}}
\newcommand{\acts}{\;\reflectbox{\rotatebox[origin=c]{-90}{$\circlearrowright$}}\;}
\newcommand{\disjunion}{\mathrel{\text{\scalebox{1.07}{$\perp\mkern-10mu\perp$}}}}
\DeclareMathOperator{\SO}{SO}
\DeclareMathOperator{\stab}{stab}
\DeclareMathOperator{\nullity}{nullity}
\DeclareMathOperator{\Perm}{Perm}
\DeclareMathOperator{\nsubgp}{\vartriangleleft}
\DeclareMathOperator{\notnsubgp}{\ntriangleleft}
\newcommand{\presentation}[2]{\left<#1\;\middle|\;#2\right>}
\DeclareMathOperator{\Char}{char}
\DeclareMathOperator{\fchar}{char}
\DeclareMathOperator{\triv}{triv}
\DeclareMathOperator{\reg}{reg}
\DeclareMathOperator{\std}{std}
\DeclareMathOperator{\Func}{Func}
\DeclareMathOperator{\End}{End}

% Convex Optimization
\let\sectionSymbol\S
\let\S\relax
\newcommand{\S}{\mathbb{S}}
\DeclareMathOperator{\dist}{dist}
\DeclareMathOperator{\dom}{dom}
\DeclareMathOperator{\diag}{diag}
\DeclareMathOperator{\ones}{\mathbbm{1}}
\newcommand{\minimizeOver}[3]{\begin{array}{rl}\underset{#1}{\ptxt{minimize}} & #2\\ \ptxt{subject to} & #3\end{array}}
\newcommand{\maximizeOver}[3]{\begin{array}{rl}\underset{#1}{\ptxt{maximize}} & #2\\ \ptxt{subject to} & #3\end{array}}
\newcommand{\minimizationProblem}[2]{\minimizeOver{}{#1}{#2}}
\newcommand{\maximizationProblem}[2]{\maximizeOver{}{#1}{#2}}
\newcommand{\minimizeOverUnconstrained}[2]{\begin{array}{rl}\underset{#1}{\ptxt{minimize}} & #2\end{array}}
\newcommand{\maximizeOverUnconstrained}[2]{\begin{array}{rl}\underset{#1}{\ptxt{maximize}} & #2\end{array}}
\newcommand{\minimizationUnconstrained}[1]{\minimizeOverUnconstrained{}{#1}}
\newcommand{\maximizationUnconstrained}[1]{\maximizeOverUnconstrained{}{#1}}
\DeclareMathOperator{\argmin}{argmin}
\DeclareMathOperator{\argmax}{argmax}

% Proofs
\newcommand{\st}{s.t.}
\newcommand{\unique}{!}
\newcommand{\iffdef}{\overset{\ptxt{def}}{\Leftrightarrow}}
\newcommand{\eqVertical}{\rotatebox[origin=c]{90}{=}}
\newcommand{\mapsfrom}{\mathrel{\reflectbox{\ensuremath{\mapsto}}}}
\newcommand{\mapsdown}{\rotatebox[origin=c]{-90}{$\mapsto$}\mkern2mu}
\newcommand{\mapsup}{\rotatebox[origin=c]{90}{$\mapsto$}\mkern2mu}
\newcommand{\from}{\!\mathrel{\reflectbox{\ensuremath{\to}}}}
\newcommand{\labeledeq}[1]{\overset{\mathclap{\ptxt{#1}}}{=}}
\newcommand{\eqdef}{\labeledeq{def}}
\newcommand{\leftsquigarrow}{\mathrel{\reflectbox{\ensuremath{\rightsquigarrow}}}}

% Brackets
\newcommand{\paren}[1]{\left(#1\right)}
\renewcommand{\brack}[1]{\left[#1\right]}
\renewcommand{\brace}[1]{\left\{#1\right\}}
\newcommand{\ang}[1]{\left<#1\right>}

% Algorithms
\algrenewcommand{\algorithmiccomment}[1]{\hskip 1em \texttt{// #1}}
\algrenewcommand\algorithmicrequire{\textbf{Input:}}
\algrenewcommand\algorithmicensure{\textbf{Output:}}
\newcommand{\algP}{\ptxt{\textbf{P}}}
\newcommand{\algNP}{\ptxt{\textbf{NP}}}
\newcommand{\algNPC}{\ptxt{\textbf{NP-Complete}}}
\newcommand{\algNPH}{\ptxt{\textbf{NP-Hard}}}
\newcommand{\algEXP}{\ptxt{\textbf{EXP}}}
\DeclareMathOperator{\fl}{fl}

%%%%%%%%%%%%%%%%%%%%%%%%%%%%%
% Other commands and shorthand
%%%%%%%%%%%%%%%%%%%%%%%%%%%%%
\newcommand{\flag}[1]{\textbf{\textcolor{red}{#1}}}
\let\uSym\u
\let\u\relax
\newcommand{\u}[1]{\underline{#1}}
\let\bSym\b
\let\b\relax
\newcommand{\b}[1]{\textbf{#1}}
\let\iSym\i
\let\i\relax
\newcommand{\i}[1]{\textit{#1}}
\let\scSym\sc
\let\sc\relax
\newcommand{\sc}[1]{\textsc{#1}}
\newcommand{\mf}[1]{\mathfrak{#1}}
\newcommand{\mc}[1]{\mathcal{#1}}
\newcommand{\ms}[1]{\mathscr{#1}}

%%%%%%%%%%%%%%%%%%%%%%%%%%%%%%%%%%%%%%%
% Make l's curvy in math environments %
%%%%%%%%%%%%%%%%%%%%%%%%%%%%%%%%%%%%%%%
\mathcode`l="8000
\begingroup
\makeatletter
\lccode`\~=`\l
\DeclareMathSymbol{\lsb@l}{\mathalpha}{letters}{`l}
\lowercase{\gdef~{\ifnum\the\mathgroup=\m@ne \ell \else \lsb@l \fi}}%
\endgroup

%%%%%%%%%%%%%%%%%%%%%%%%%
% Fix \vdots and \ddots %
%%%%%%%%%%%%%%%%%%%%%%%%%
\usepackage{letltxmacro}
\LetLtxMacro\orgvdots\vdots
\LetLtxMacro\orgddots\ddots

\makeatletter
\DeclareRobustCommand\vdots{%
	\mathpalette\@vdots{}%
}
\newcommand*{\@vdots}[2]{%
	% #1: math style
	% #2: unused
	\sbox0{$#1\cdotp\cdotp\cdotp\m@th$}%
	\sbox2{$#1.\m@th$}%
	\vbox{%
		\dimen@=\wd0 %
		\advance\dimen@ -3\ht2 %
		\kern.5\dimen@
		% remove side bearings
		\dimen@=\wd2 %
		\advance\dimen@ -\ht2 %
		\dimen2=\wd0 %
		\advance\dimen2 -\dimen@
		\vbox to \dimen2{%
			\offinterlineskip
			\copy2 \vfill\copy2 \vfill\copy2 %
		}%
	}%
}
\DeclareRobustCommand\ddots{%
	\mathinner{%
		\mathpalette\@ddots{}%
		\mkern\thinmuskip
	}%
}
\newcommand*{\@ddots}[2]{%
	% #1: math style
	% #2: unused
	\sbox0{$#1\cdotp\cdotp\cdotp\m@th$}%
	\sbox2{$#1.\m@th$}%
	\vbox{%
		\dimen@=\wd0 %
		\advance\dimen@ -3\ht2 %
		\kern.5\dimen@
		% remove side bearings
		\dimen@=\wd2 %
		\advance\dimen@ -\ht2 %
		\dimen2=\wd0 %
		\advance\dimen2 -\dimen@
		\vbox to \dimen2{%
			\offinterlineskip
			\hbox{$#1\mathpunct{.}\m@th$}%
			\vfill
			\hbox{$#1\mathpunct{\kern\wd2}\mathpunct{.}\m@th$}%
			\vfill
			\hbox{$#1\mathpunct{\kern\wd2}\mathpunct{\kern\wd2}\mathpunct{.}\m@th$}%
		}%
	}%
}
\makeatother

\newcommand{\B}{
	\begin{tikzpicture}
	\filldraw [fill=red, draw=black] (0, 0) rectangle (0.37, 0.45);
	\draw [line width=0.5mm, white ] (0.1,0.08) -- (0.1,0.38);
	\draw[line width=0.5mm, white ] (0.1, 0.35) .. controls (0.2, 0.35) and (0.4, 0.2625) .. (0.1, 0.225);
	\draw[line width=0.5mm, white ] (0.1, 0.225) .. controls (0.2, 0.225) and (0.4, 0.1625) .. (0.1, 0.1);
	\end{tikzpicture}
}

% Allow custom symbols as arrows in tikz-cd
\tikzset{
  symbol/.style={
    draw=none,
    every to/.append style={
      edge node={node [sloped, allow upside down, auto=false]{$#1$}}}
  }
}
% Example: \arrow[r,symbol=\cong]
% https://tex.stackexchange.com/questions/394154/how-to-include-inclusion-subgroup-relationship-in-tikz-cd-diagram

\author{Thomas Cohn}
\title{Math 635 Lecture 36}
\date{4/14/21} % Can also use \today

\begin{document}
\maketitle
\setlength\RaggedRightParindent{\parindent}
\RaggedRight

\par\noindent
Recall from last time: Let $M$ be a Riemannian manifold. We defined the differential operators $\del$ (gradient) and $\div$ (divergence), and we have
\[
	\begin{tikzcd}
		C^{\infty}(M) \arrow[r,shift left, "\del"] \arrow[d, symbol={=}] & \mf{X}(M) \arrow[l,shift left, "\scalebox{0.75}{$\begin{smallmatrix}\div\\ \ptxt{or}\\ -\div\end{smallmatrix}$}"] \arrow[d,symbol=\cong,"\,\,\ptxt{(metric dual)}"]\\
		C^{\infty}(M) \arrow[r, shift left, "d"] & \Omega^{1}(M) \arrow[l, shift left, "\scalebox{0.75}{$\begin{smallmatrix}-\delta\\ \ptxt{or}\\ \delta\end{smallmatrix}$}"]
	\end{tikzcd}
\]
``$\delta=-\div$ on the differential form side''\n

\par\noindent
We also defined the Laplacian on functions $\Laplacian:C^{\infty}(M)\to{}C^{\infty}(M)$ by $\Laplacian=\delta\of{}d$ iff $\Laplacian{}f=-\div(\del{}f)$.\n

\subsection*{$l^{2}$ Inner Products}

\defn{
	Assume $M$ is oriented. $\forall{}f,g\in{}C^{\infty}(M)$, we define the \u{$l^{2}$ inner product} by
	\[
		\iprod{f,g}_{l^{2}}=\int_{M}fg\,\,d\Vol
	\]
}

\par\noindent
We can extend this to sections of real vector bundles over $M$, $\mc{E}\overset{\pi}{\to}M$. Put a Euclidean structure on the fibers of $\mc{E}$: $\forall{}p\in{}M$, $\giprod_{p}$ is a Euclidean inner product on $\mc{E}_{p}=\pi\inv(p)$, varying smoothly with $p$.\n

\defn{
	$\forall{}s,t\in\Gamma_{0}(\mc{E})$ (compactly supported sections). Then we define the \u{$l^{2}$ inner product} by
	\[
		\iprod{s,t}_{l^{2}}=\int_{M}\underbrace{\iprod{s(p),t(p)}_{p}}_{\ptxt{function of $p$}}d\Vol
	\]
}

\par\noindent
Consider the case $\mc{E}=\bigwedge^{k}(T^{*}M)$. Then the Euclidean structure on $\bigwedge^{k}(T^{*}M)$ is induced by the Riemannian metric. For $k=1$, we simply have $T^{*}M\cong{}TM$ by the metric dual. For general $k$, $\forall{}p\in{}M$, let $V=T_{p}^{*}M$. We define an inner product on $\bigwedge^{k}V$:
\[
	(v_{1}\wedge\cdots\wedge{}v_{k},w_{1}\wedge\cdots\wedge{}w_{k})\eqdef\det\paren{\iprod{v_{i},w_{j}}}_{ij}
\]
Check: If $(e_{1},\ldots,e_{n})$ is an orthonormal basis of $V$, then $\set{e_{i_{1}}\wedge\cdots\wedge{}e_{i_{k}}\mid{}i_{1}<\cdots<i_{k}}$ is an orthonormal basis of $\bigwedge^{k}V$.\n

\par\noindent
In this way, we get the notion of an $l^{2}$ inner product of any two $k$-forms $\alpha,\beta\in\Omega^{k}(M)$ by
\[
	\iprod{\alpha,\beta}_{l^{2}}=\int_{M}\iprod{\alpha_{p},\beta_{p}}_{p}d\Vol
\]

\prop{
	$\forall{}f\in{}C^{\infty}(M),X\in\mf{X}(M)$, one has
	\[
		\iprod{\del{}f,X}_{l^{2}}=-\iprod{f,\div{}X}_{l^{2}}
	\]
	That is, $\forall{}f\in\Omega^{0}(M),\alpha\in\Omega^{1}(M)$,
	\[
		\iprod{df,\alpha}_{l^{2}}=\iprod{f,\delta\alpha}_{l^{2}}
	\]
	That is, $\delta=d^{*}$, the adjoint of $d$, so $\Laplacian=d^{*}d$.\nn
	Proof: Start with $\mc{L}_{fX}(d\Vol)=f\mc{L}_{X}(d\Vol)+Xf$. Now integrate:
	\[
		\int_{M}\mc{L}_{fX}(d\Vol)=\int_{M}\div(fX)d\Vol=0
	\]
	because $\d{}M=\emptyset$. So we have
	\[
		0=\int_{M}f\div(X)d\Vol+\int_{M}\underbrace{\iprod{X,\del{}F}}_{\hspace{-35pt}\mathrlap{X(f)=df(X)=\iprod{\del{}f,X}}}d\Vol
	\]
	So $0=\iprod{f,\div{}X}_{l^{2}}+\iprod{X,\del{}f}_{l^{2}}$.\proven
}

\cor{
	$\iprod{\Laplacian{}f,g}_{l^{2}}=\iprod{f,\Laplacian{}g}_{l^{2}}$.\n
}

\par\noindent
Now, generalize to $\Omega^{k}$. (The previous discussion was for $k=0$.)
\[
	\begin{tikzcd}
		\Omega^{k} \arrow[r, shift left, "d"] & \Omega^{k+1} \arrow[l, shift left, "\delta=d^{*}=?"]
	\end{tikzcd}
\]
Is there a $\delta$? What is it?\n

\par\noindent
In local coordinates, $\delta$ is \i{also} a differential operator of degree $1$. Try integrating by parts!\n

\par\noindent
Preliminary linear algebra: the Hodge star operator. Let $V$ be an $n$-dimensional vector space, oriented, with an inner product. We claim that $\forall{}k$, there exists $\star:\bigwedge^{k}V\to\bigwedge^{n-k}V$ linear such that for any positive oriented basis $(e_{1},\ldots,e_{n})$ of $V$, $\star(e_{1}\wedge\cdots\wedge{}e_{k})=e_{k+1}\wedge\cdots\wedge{}e_{n}$.\n

\ex{
	For $V=\R^{3}$ with the standard orientation,
	\[
		\map{\star:\bigwedge^{2}V}{\bigwedge^{1}V}{dx^{1}\wedge{}dx^{2}}{dx^{3}}
	\]
	(Now do it cyclically.)\n
}

\par\noindent
Note: $\dim\bigwedge^{k}=\binom{n}{k}=\binom{n}{n-k}=\dim\bigwedge^{n-k}$.\n

\par\noindent
Observe: On $\R^{3}$ in the calc 3 context, for $X\in\mf{X}(\R^{3})$, we define
\[
	\curl{}X=\del\times{}X\in\mf{X}(M)
\]
What is this object? Well,
\[
	\begin{tikzcd}
		\mf{X}(\R^{3}) \arrow[r,symbol=\cong] &[-15pt] \Omega^{1}(\R^{3}) \arrow[r,"d"] \arrow[r, bend right=30, ""{coordinate, name=Z1}, phantom] & \Omega^{2}(\R^{3}) \arrow[r,"\star"] \arrow[r, bend right=30, ""{coordinate, name=Z2}, phantom] & \Omega^{1}(\R^{3}) \arrow[r,symbol=\cong] &[-15pt] \mf{X}(\R^{3})
		\arrow[from=1-1, to=1-5, to path={
			-- (\tikztostart.south)
			|- (Z1)
			-- (Z2)\tikztonodes
			-| (\tikztotarget.south)
			-- (\tikztotarget.south)
		}, "\curl"']
	\end{tikzcd}
\]
Note that this only works for $\dim{}=3$.\n

\par\noindent
Some properties of $\star$:
\begin{enumerate}
	\item We have
	\[
		\begin{tikzcd}
			\bigwedge^{k} \arrow[r,"\star"] \arrow[r, bend right=45, ""{coordinate, name=Z1}, phantom] & \bigwedge^{n-k} \arrow[r,"\star"] \arrow[r,bend right=45, ""{coordinate, name=Z2}, phantom] & \bigwedge^{k}
			\arrow[from=1-1, to=1-3, to path={
				-- (\tikztostart.south)
				|- (Z1)
				-- (Z2)\tikztonodes
				-| (\tikztotarget.south)
				-- (\tikztotarget.south)
			}, "(-1)^{k(n-k)}\Id"']
		\end{tikzcd}
	\]
	because
	\[
		e_{1}\wedge\cdots\wedge{}e_{k}\overset{\star}{\mapsto}e_{k+1}\wedge\cdots\wedge{}e_{n}\overset{\star}{\mapsto}(-1)^{\sigma}e_{1}\wedge\cdots\wedge{}e_{k}
	\]
	``$n-k$ signs, $k$ times''.
	\item $\star:\bigwedge^{n}V\to\bigwedge^{0}V=\R$ has $\star(\Vol)=1$.
	\item $\forall\alpha,\beta\in\bigwedge^{k}V$, $\iprod{\alpha,\beta}=\star(\alpha\wedge(\star\beta))\in\R$.
\end{enumerate}

\cor{
	Apply/extend $\star$ to forms on a compact, oriented, Riemannian manifold $M$ (with $\dim{}M=n$), $\Omega^{k}(M)$, by acting pointwise: $\star:\Omega^{k}(M)\to\Omega^{n-k}(M)$. Note: $\forall\alpha,\beta\in\Omega^{k}(M)$, $\iprod{\alpha,\beta}_{l^{2}}=\int_{M}\alpha\wedge(\star\beta)$.\n
}

\par\noindent
Back to our main question:\n

\prop{
	The adjoint $\delta$ of $d:\Omega^{k}\to\Omega^{k+1}$ is $\delta=(-1)^{nk+1}\star{}d\star$.\n
}

\par\noindent
Note: If $\beta\in\Omega^{k+1}$, $\star\beta\in\Omega^{n-k-1}$, so $d\star\beta\in\Omega^{n-k}$, so $\star{}d\star\beta\in\Omega^{k}$. Superficially, $\delta=\star{}d\star:\Omega^{k+1}\to\Omega^{k}$. Now, we prove it:\n

\par\noindent
Proof: Let $\alpha\in\Omega^{k},\beta\in\Omega^{k+1}$. We want to show $\iprod{d\alpha,\beta}_{l^{2}}=\iprod{\alpha,\delta\beta}_{l^{2}}$. We'll use integration by parts. Starting with the fact that $0=\int_{M}d(\alpha\wedge\star\beta)$, because $\alpha\wedge\star\beta$ is a $n-1$ form, so $d(\alpha\wedge\star\beta)$ is a top-degree form. By Stokes' theorem, since we have an empty boundary, this integral is $0$. Well,
\[
	0=\int_{M}d(\alpha\wedge\star\beta)=\underbrace{\int_{M}d\alpha\wedge\star\beta}_{=\iprod{d\alpha,\beta}_{l^{2}}}+(-1)^{k}\int_{M}\alpha\wedge(d\star\beta)
\]
So
\[
	\iprod{d\alpha,\beta}_{l^{2}}=(-1)^{?}\int_{M}\alpha\wedge(d\star\beta)=(-1)^{?}\int_{M}\alpha\wedge(\star\star)d\star\beta=(-1)^{?}\iprod{\alpha,\star{}d\star\beta}_{l^{2}}=(-1)^{?}\iprod{\alpha,\delta\beta}
\]
(We didn't do the sign computations, but they do work out.)\proven

\defn{
	The Laplacian on forms $\Laplacian:\Omega^{k}(M)\to\Omega^{k}(M)$ is $\Laplacian=\delta{}d+d\delta$.\n
}

\end{document}